\section{Розділ~\ref{sec:motorlearning}. Навчання рухів}

Правило найменших квадратів і вигода окремого проводу помилки --- теореми; будову схеми з проводом помилки, послаблення за збігу, підлаштування сліду під затримку, післядію й спонтанне повернення навички виміряно; те, що вся схема працює як навчання з учителем і будує внутрішню модель, --- провідна гіпотеза; числа двопроцесної моделі --- розрахунок за моделлю книги.

{\footnotesize
\setlength{\tabcolsep}{3pt}\renewcommand{\arraystretch}{1.15}
\begin{longtable}{>{\raggedright}p{0.5cm}>{\raggedright}p{4.0cm}>{\raggedright}p{5.6cm}>{\raggedright}p{2.3cm}>{\raggedright\arraybackslash}p{1.7cm}}
\toprule
№ & Твердження & Дослід і що виміряно & Джерело & Статус \\
\midrule\endfirsthead
\toprule
№ & Твердження & Дослід і що виміряно & Джерело & Статус \\
\midrule\endhead
1 & Правило~\eqref{eq:lms} у середньому йде за градієнтом квадрата помилки & Результат теорії адаптивних фільтрів: поправка, пропорційна входу й помилці виходу, є стохастичним спуском за градієнтом середнього квадрата помилки. & \cite{WidrowHoff1960} & теорія \\
2 & З проводом помилки на кожен вихід швидкість не залежить від кількості виходів; з одним спільним числом --- падає з кількістю нейронів, що шумлять (твердження~\ref{prop:teachers}) & Криві навчання лінійної мережі: навчання за випадковими збуреннями нейронів повільніше за точний градієнт у стільки разів, скільки нейронів збурюють. & \cite{Werfel2005} & теорія \\
3 & Схема з проводом помилки працює як навчання з учителем (твердження~\ref{prop:errorwire}) & Теорії, виведені з будови схеми. Досліди рядків 7--10 з ними узгоджуються; що вся схема працює саме так, не доведено. & \cite{Marr1969, Albus1971} & гіпотеза \\
4 & Розширювач: близько $7\cdot10^{10}$ дрібних \nt{GLUT}-нейронів, більше, ніж у всьому іншому мозку & Підрахунок ядер у суспензії розчиненої тканини: у мозочку людини $69\cdot10^9$ нейронів із $86\cdot10^9$ в усьому мозку. Стереологія: $101\cdot10^9$ дрібних нейронів у мозочку літніх чоловіків. & \cite{Azevedo2009, Andersen1992cereb} & виміряно \\
5 & Підняття розмірності входу робить потрібну залежність лінійною & Теорема про кількість розбиттів: зважена сума $n$ входів із порогом реалізує довільну розмітку приблизно $2n$ векторів загального положення. Що мозочок користується саме цим, --- частина гіпотези рядка~3. & \cite{Cover1965} & теорія \\
6 & Близько $3\cdot10^7$ вихідних \nt{GABA}-нейронів, у кожного порядку $10^5$ входів & Стереологія мозочка людини: $30.5\cdot10^6$ вихідних нейронів. Електронна мікроскопія, щур: близько $1.75\cdot10^5$ входів від розширювача на один вихідний нейрон. Гальмівний медіатор вихідних нейронів --- в огляді. & \cite{Andersen1992cereb, NapperHarvey1988, Ito2001} & виміряно \\
7 & Рівно один провід помилки на вихідний нейрон, близько разу на секунду, щоразу потужний спалах & Огляд записів: кожен вихідний нейрон отримує один ліаноподібний \nt{GLUT}-вхід, що викликає складний імпульс із частотою порядку $1$~Гц. & \cite{Ito2001} & виміряно \\
8 & Імовірність спалаху залежить від напрямку помилки руху & Мавпа, окорухова частина мозочка, адаптація саккад: напрямок зорової помилки задає ймовірність складного імпульсу, величина --- його час; спалах змінює прості імпульси в наступній спробі й зсуває рух уздовж переважної помилки клітини. & \cite{Herzfeld2018} & виміряно \\
9 & Збіг спалаху помилки з активним входом послаблює вхід ($-\eta e_i x_j$) & Огляд дослідів: спільна стимуляція входу від розширювача й проводу помилки викликає довготривале послаблення цього входу; стимуляція одного входу --- ні. & \cite{Ito2001} & виміряно \\
10 & Довжину сліду підлаштовано під затримку помилки: за затримки $100\ms$ найсильніше послаблюється вхід, що був активним за $100\ms$ & Зрізи й живі миші: у частині, що відповідає за окорухове навчання, послаблення вузько налаштоване на інтервал близько $120\ms$ між входом і помилкою --- стільки йде сигнал помилки в цій задачі; в іншій частині різні клітини налаштовано на різні інтервали. & \cite{Suvrathan2016} & виміряно \\
11 & Схема --- адаптивний фільтр~\eqref{eq:adaptivefilter}, що вивчає внутрішню модель тіла & Модель, виведена з будови схеми. Прямого досліду, де вивчений фільтр виміряно як передавальну функцію тіла, немає. & \cite{Fujita1982} & гіпотеза \\
12 & Передбачення віднімає наслідки власних рухів, тому себе не полоскочеш & Люди: власний дотик здається менш лоскітним, ніж такий самий чужий; у томографі соматосенсорна кора відповідає на власний дотик слабше, а мозочок менш активний, коли рух торкається шкіри. Пояснення через віднімання передбачення --- гіпотеза. & \cite{Blakemore1998} & гіпотеза \\
13 & За зовнішньої сили промах зменшується, а після її зняття рука промахується в інший бік & Люди тягнуться рукояттю робота в силовому полі: траєкторії повертаються до прямих; після зняття поля післядія --- майже дзеркальне відображення перших викривлень; вона переноситься й у нетреновані частини простору. & \cite{ShadmehrMussaIvaldi1994} & виміряно \\
14 & Навчаються два процеси~\eqref{eq:tworate}; після зворотного збурення навичка повертається сама & Люди, силове поле, далі коротке зворотне поле й спроби із затиснутою помилкою: поправка сама повертається до старої; модель із двох процесів це відтворює, модель з одного --- ні. & \cite{Smith2006} & виміряно \\
15 & Числа рис.~\ref{fig:tworate}: $0.84$ (повільний $0.63$) за $200$ рухів; $6$ зворотних; швидкий $-0.53$, повільний $0.46$; повернення до $0.37$ за $40$ рухів & Розрахунок за \texttt{models/\allowbreak motor\_\allowbreak learning.py}, $A_f=0.92$, $B_f=0.1$, $A_s=0.996$, $B_s=0.02$: $0.840$ ($0.634$ і $0.205$), $6$ рухів, $-0.531$ і $0.457$, пік $0.371$ на $40$-му русі. & \texttt{models/\allowbreak motor\_\allowbreak learning.py} & модель \\
16 & Два процеси потрібні, бо тіло змінюється з різними швидкостями & Теорія: оптимальна оцінка за завад із кількома часами життя дає кілька фільтрів із різними сталими часу й пояснює спонтанне повернення та прискорене повторне навчання. & \cite{Kording2007} & гіпотеза \\
\bottomrule
\end{longtable}}
